\subsection{Desafíos en la enseñanza de la trigonometría}

\subsubsection{Limitaciones del enfoque curricular tradicional}

La trigonometría ocupa un lugar central en el currículo de matemáticas de educación media, constituyendo un puente fundamental entre la geometría euclidiana y el análisis matemático avanzado \parencite{weberConnectingResearchTeaching2008} Los lineamientos curriculares del Ministerio de Educación Nacional reconocen este campo como esencial para el desarrollo del pensamiento variacional y funcional \parencite{ministeriodeeducacionnacionalSerieLineamientosCurriculares1988}. No obstante, diversos estudios revelan que su implementación en el aula presenta serias deficiencias.

 Mayormente en la educación media, los contenidos trigonométricos se circunscriben casi exclusivamente al estudio de las razones del triángulo rectángulo: seno, coseno y tangente \parencite{fiallolealAcercaInvestigacionEducacion2015}. Esta aproximación limitada omite elementos centrales del pensamiento trigonométrico, tales como la interpretación funcional de estas razones, el análisis sistemático de periodicidad o la introducción al estudio de funciones trascendentes.

Una consecuencia de esta omisión curricular es no  vincular la trigonometría con procesos de modelización en contextos reales \parencite{rosaEnsenanzaAprendizajeTrigonometria2023}. Los estudiantes desarrollan una visión instrumental del conocimiento trigonométrico, percibiéndolo únicamente como un conjunto de fórmulas para resolver problemas geométricos específicos \parencite{delgadoEstrategiaDidacticaPara2023}. Esta perspectiva resulta insuficiente cuando los educandos avanzan a cursos superiores donde deben aplicar conceptos trigonométricos para modelar fenómenos periódicos en física, analizar señales en ingeniería o comprender procesos oscilatorios en ciencias naturales \parencite{orhaniAddressingStudentsChallenges2024}.

\subsubsection{Dificultades conceptuales persistentes}

Diversas investigaciones coinciden en señalar que los estudiantes presentan obstáculos conceptuales persistentes tanto en la comprensión como en la aplicación de funciones trigonométricas . Estas dificultades trascienden las meras deficiencias procedimentales y revelan problemas más profundos en la construcción del conocimiento trigonométrico. Entre las dificultades más reportadas se encuentran: la confusión sistemática entre razones trigonométricas y relaciones métricas; una dependencia excesiva de procedimientos algorítmicos que impide la comprensión conceptual; escasa interpretación de las funciones trigonométricas en términos de variación y cambio; y una comprensión limitada del uso de la medida en radianes y del comportamiento periódico de estas funciones \parencite{delgadoEstrategiaDidacticaPara2023,iriarteDificultadesAprendizajeEcuaciones2025,chaconComprensionRazonesTrigonometricas2007}.

Estas dificultades no son meramente académicas. Se traducen en bajos niveles de rendimiento en evaluaciones estandarizadas y, más preocupante aún, en una limitada apropiación de las aplicaciones de la trigonometría en áreas disciplinares relacionadas como la física, la acústica o diversas ramas de la ingeniería \parencite{andradeSuperficialidadEnsenanzaTrigonometria2020}. Los estudiantes que no logran superar estos obstáculos conceptuales podrían enfrentar limitaciones para avanzar en su formación académica superior.

\subsubsection{Predominio del enfoque geométrico-algorítmico}

Otro elemento crítico identificado en la literatura es la persistencia de un enfoque eminentemente geométrico y algorítmico en la enseñanza de la trigonometría \parencite{magatFunctionalApproachTeaching2018,rosaEnsenanzaAprendizajeTrigonometria2023}. Esta aproximación pedagógica presenta la trigonometría exclusivamente como un conjunto de relaciones entre los lados y ángulos de un triángulo rectángulo, restringiendo dramáticamente su alcance conceptual.

Este tratamiento tradicional dificulta que los estudiantes reconozcan las funciones trigonométricas como entidades analíticas, periódicas y modelizadoras de fenómenos naturales y tecnológicos. El énfasis desproporcionado en la memorización de fórmulas y en la ejecución mecánica de procedimientos ha contribuido a una aproximación superficial y descontextualizada del conocimiento trigonométrico 
\parencite{nguDifferentialInstructionalEffectiveness2023}.

Como resultado de este enfoque limitado, los estudiantes encuentran serias dificultades para transferir lo aprendido a nuevos contextos, para establecer conexiones significativas entre la trigonometría y otras áreas del conocimiento matemático, y para reconocer la utilidad práctica de estos conceptos en la resolución de problemas auténticos. La trigonometría se percibe como un contenido aislado, carente de conexiones con la experiencia cotidiana o con aplicaciones significativas.

\subsubsection{Necesidad de transformación pedagógica}
Otro elemento crítico identificado en la literatura es la persistencia de un enfoque eminentemente geométrico y algorítmico en la enseñanza de la trigonometría, lo que ha sido señalado como una de las principales limitaciones de su implementación curricular. Diversos estudios muestran que este tratamiento se centra casi exclusivamente en el uso del triángulo rectángulo y en la aplicación mecánica de razones trigonométricas, dejando de lado la comprensión profunda de la función trigonométrica y su papel en la modelación de fenómenos periódicos. Un ejemplo de este señalamiento puede encontrarse en el trabajo de \textcite{chavezEstrategiasDidacticasMediadas2022}, quienes evidencian cómo la enseñanza tradicional enfatiza la resolución de ejercicios rutinarios y la manipulación algorítmica,


La literatura sugiere que superar las limitaciones actuales requiere ir más allá de simples ajustes metodológicos. Se necesita una transformación profunda que incluya: la reconceptualización del currículo para enfatizar una visión más amplia del pensamiento trigonométrico; el desarrollo de estrategias pedagógicas que promuevan la comprensión conceptual sobre la memorización mecánica; la integración sistemática de contextos reales que demuestren la utilidad de la trigonometría; y la incorporación de tecnologías que faciliten la exploración dinámica de conceptos y la personalización del aprendizaje.
